% ── Answer Key: Practice 1 — Set Theory ──────────────────────────
% This file is meant to be \input{} into a master document.
% It assumes amsmath, amssymb, amsthm are loaded and
% \newcommand{\comp}[1]{\overline{#1}}, \newcommand{\symdiff}{\mathbin{\triangle}} exist.

\section*{Գործնական 1 — Բազմությունների տեսություն}

%======================================================================
\subsection*{Խնդիր 1}
Ի՞նչ է ներկայացնում $\{\varnothing\}$ բազմությունը և քանի՞ տարր ունի այն։

$\{\varnothing\}$-ն բազմություն է, որի միակ տարրը դատարկ բազմությունն է։ Այն ունի $\boxed{1}$ տարր։

%======================================================================
\subsection*{Խնդիր 2}
Ո՞ր պնդումներն են ճիշտ։

\begin{itemize}
  \item[2.1.] $\varnothing \in \{\varnothing, \{\varnothing\}\}$ — \textbf{ՃԻՇՏ Է}։ Բազմությունը պարունակում է $\varnothing$-ն որպես տարր։
  \item[2.2.] $\{\varnothing\} \in \{\{\varnothing\}\}$ — \textbf{ՃԻՇՏ Է}։ $\{\{\varnothing\}\}$ բազմության միակ տարրը $\{\varnothing\}$-ն է։
  \item[2.3.] $\varnothing \in \{\{\varnothing\}\}$ — \textbf{ՍԽԱԼ Է}։ Միակ տարրը $\{\varnothing\}$-ն է, որը հավասար չէ $\varnothing$-ին։
\end{itemize}

%======================================================================
\subsection*{Խնդիր 3}
Քանի՞ տարր ունի $\{\{\varnothing\}\}$ բազմությունը։

$\boxed{1}$։ Նրա միակ տարրը $\{\varnothing\}$-ն է։

%======================================================================
\subsection*{Խնդիր 4}
Քանի՞ տարր ունեն հետևյալ բազմությունները։

\begin{itemize}
  \item[4.1.] $\{1,2,1\}$՝ $\boxed{2}$ (կրկնությունները անտեսվում են, բազմությունը $\{1,2\}$-ն է)։
  \item[4.2.] $\{1,2,\{1,2\}\}$՝ $\boxed{3}$ (տարրերն են՝ $1$, $2$, $\{1,2\}$)։
  \item[4.3.] $\{\{2\}\}$՝ $\boxed{1}$ (տարրն է՝ $\{2\}$)։
  \item[4.4.] $\{1,\{1\}\}$՝ $\boxed{2}$ (տարրերն են՝ $1$ և $\{1\}$, դրանք տարբեր են)։
  \item[4.5.] $\{1,\varnothing\}$՝ $\boxed{2}$ (տարրերն են՝ $1$ և $\varnothing$)։
  \item[4.6.] $\{\{\varnothing\},\varnothing\}$՝ $\boxed{2}$ (տարրերն են՝ $\{\varnothing\}$ և $\varnothing$)։
  \item[4.7.] $\{\{\varnothing\},\{\varnothing\}\}$՝ $\boxed{1}$ (կրկնություն, բազմությունը $\{\{\varnothing\}\}$-ն է)։
  \item[4.8.] $\{x \mid 2x=1\}$, որտեղ $x \in \mathbb{R}$՝ $\boxed{1}$ (բազմությունը $\{\tfrac{1}{2}\}$-ն է)։
\end{itemize}

%======================================================================
\subsection*{Խնդիր 5}
Ճի՞շտ է արդյոք, որ $\{1,2\} \in \{1,2,\{1,3\},\{1,2,3\}\}$։

\textbf{ՈՉ։} Աջ կողմի բազմության տարրերն են՝ $1$, $2$, $\{1,3\}$ և $\{1,2,3\}$։ Դրանցից ոչ մեկը հավասար չէ $\{1,2\}$-ին։

%======================================================================
\subsection*{Խնդիր 6}
Բերեք բազմությունների օրինակներ, որոնք իրենք իրենց տարր են։

ZFC բազմությունների ստանդարտ տեսությունում (Ռեգուլյարության աքսիոմով) ոչ մի բազմություն իր տարրը չէ։ Սակայն, «նաիվ» տեսությունում կամ առանց ռեգուլյարության տեսություններում, «բոլոր բազմությունների բազմությունը»՝ $V = \{x \mid x = x\}$, կբավարարեր $V \in V$ պայմանին։ Մեկ այլ դասական օրինակ՝ եթե $S$-ը սահմանենք որպես «բոլոր այն բազմությունների բազմություն, որոնք թվերի բազմություններ չեն», ապա $S \in S$ (Ռասելյան տիպի կառուցում)։ ZFC-ում \textbf{նման բազմություն գոյություն չունի}։

%======================================================================
\subsection*{Խնդիր 7}
Բերեք $A, B, C$ բազմությունների օրինակներ, որտեղ $A \in B$, $B \in C$, բայց $A \notin C$։

Դիցուք $A = 1$, $B = \{1,2\}$, $C = \{\{1,2\},3\}$։
Այդ դեպքում $1 \in \{1,2\}$, $\{1,2\} \in \{\{1,2\},3\}$, բայց $1 \notin \{\{1,2\},3\}$, քանի որ $C$-ի տարրերն են $\{1,2\}$-ը և $3$-ը։

%======================================================================
\subsection*{Խնդիր 8}
Ապացուցել՝ $\{\{a\},\{a,b\}\} = \{\{c\},\{c,d\}\} \iff a=c$ և $b=d$։

\textit{Ապացույց.}
($\Leftarrow$) Ակնհայտ է տեղադրմամբ։
($\Rightarrow$) Հավասարությունից հետևում է, որ $\{a\}$-ն պետք է հավասար լինի կա՛մ $\{c\}$-ին, կա՛մ $\{c,d\}$-ին։
Եթե $\{a\} = \{c,d\}$, ապա $c = d = a$, և այդ դեպքում $\{a,b\} = \{c\} = \{a\}$, ինչից հետևում է $b = a = c = d$։
Եթե $\{a\} = \{c\}$, ապա $a = c$, և մնացած տարրը տալիս է $\{a,b\} = \{c,d\} = \{a,d\}$, հետևաբար $b = d$։ \qed

%======================================================================
\subsection*{Խնդիր 9}
Ապացուցել երեք բազմային հավասարություններ/ներառումներ։

\textbf{(a)} $\{x \in \mathbb{Z} \mid \exists y,\, x = 6y\} = \{x \in \mathbb{Z} \mid \exists u,v,\, x = 2u \text{ և } x = 3v\}$։

\textit{Ապացույց.} ($\subseteq$) Եթե $x = 6y$, ապա $x = 2(3y)$ և $x = 3(2y)$։
($\supseteq$) Եթե $x = 2u$ և $x = 3v$, ապա $x$-ը բաժանվում է և՛ 2-ի, և՛ 3-ի, հետևաբար նաև $\text{ԱԸԲ}(2,3) = 6$-ի։ \qed

\textbf{(b)} $\{x \in \mathbb{R} \mid \exists y,\, x = y^2\} = \{x \in \mathbb{R} \mid x \ge 0\}$։

\textit{Ապացույց.} ($\subseteq$) $x = y^2 \ge 0$։
($\supseteq$) Եթե $x \ge 0$, վերցնենք $y = \sqrt{x}$, այդ դեպքում $x = y^2$։ \qed

\textbf{(c)} $\{x \in \mathbb{Z} \mid \exists y,\, x = 6y\} \subset \{x \in \mathbb{Z} \mid \exists y,\, x = 2y\}$։

\textit{Ապացույց.} Եթե $x = 6y$, ապա $x = 2(3y)$, ուստի $\subseteq$ տեղի ունի։ Ներառումը խիստ է, քանի որ $2 = 2 \cdot 1$-ը աջ կողմի բազմությունում է, բայց $2$-ը 6-ի բազմապատիկ չէ։ \qed

%======================================================================
\subsection*{Խնդիր 10}
Ապացուցել $\subseteq$ և $\subset$ հարաբերությունների տրանզիտիվությունը։

\textbf{(a)} Եթե $A \subseteq B$ և $B \subseteq C$, ապա $A \subseteq C$։

\textit{Ապացույց.} Դիցուք $x \in A$։ Այդ դեպքում $x \in B$ (քանի որ $A \subseteq B$), որտեղից էլ $x \in C$ (քանի որ $B \subseteq C$)։ \qed

\textbf{(b)} Եթե $A \subseteq B$ և $B \subset C$, ապա $A \subset C$։

\textit{Ապացույց.} Ըստ (a)-ի, $A \subseteq C$։ Քանի որ $B \subset C$, գոյություն ունի $z \in C \setminus B$։ Եթե $A = C$, ապա $z \in A \subseteq B$, ինչը հակասություն է։ Ուստի $A \neq C$, այսինքն՝ $A \subset C$։ \qed

\textbf{(c)} Եթե $A \subset B$ և $B \subseteq C$, ապա $A \subset C$։

\textit{Ապացույց.} Ըստ (a)-ի, $A \subseteq C$։ Քանի որ $A \subset B$, գոյություն ունի $z \in B \setminus A$, և $z \in C$, քանի որ $B \subseteq C$։ Այսպիսով $z \in C \setminus A$, հետևաբար $A \neq C$։ \qed

\textbf{(d)} Եթե $A \subset B$ և $B \subset C$, ապա $A \subset C$։

\textit{Ապացույց.} Հետևում է (c)-ից, քանի որ $B \subset C$ ենթադրում է $B \subseteq C$։ \qed

%======================================================================
\subsection*{Խնդիր 11}
Բերեք $A, B, C, D, E$ բազմությունների օրինակներ, որտեղ $A \subseteq B$, $B \in C$, $C \subseteq D$, $D \subseteq E$։

$A = \{1\}$, $B = \{1,2\}$, $C = \{\{1,2\},3\}$, $D = \{\{1,2\},3,4\}$, $E = \{\{1,2\},3,4,5\}$։

%======================================================================
\subsection*{Խնդիր 12}
Ո՞ր պնդումներն են ճիշտ կամայական $A, B, C$ բազմությունների համար։

\begin{itemize}
  \item[(a)] Եթե $A \notin B$ և $B \notin C$, ապա $A \notin C$։ — \textbf{ՍԽԱԼ Է}։
    Հակաօրինակ՝ $A = 1$, $B = \{2\}$, $C = \{1\}$։ Այդ դեպքում $1 \notin \{2\}$ և $\{2\} \notin \{1\}$, բայց $1 \in \{1\}$։
  \item[(b)] Եթե $A \neq B$ և $B \neq C$, ապա $A \neq C$։ — \textbf{ՍԽԱԼ Է}։
    Հակաօրինակ՝ $A = C = \{1\}$, $B = \{2\}$։
  \item[(c)] Եթե $A \in B$ և $B \subseteq C$, ապա $A \in C$։ — \textbf{ՃԻՇՏ Է}։
    Քանի որ $A \in B$ և $B$-ի յուրաքանչյուր տարր $C$-ում է, ստանում ենք $A \in C$։
  \item[(d)] Եթե $A \subseteq B$ և $B \subseteq C$, ապա $C \subseteq A$։ — \textbf{ՍԽԱԼ Է}։
    Հակաօրինակ՝ $A = \varnothing$, $B = C = \{1\}$։
  \item[(e)] Եթե $A \subseteq B$ և $B \in C$, արդյո՞ք $A \notin C$։ — \textbf{ՍԽԱԼ Է} (պարտադիր չէ)։
    Հակաօրինակ՝ $A = \{1\}$, $B = \{1,2\}$, $C = \{\{1\},\{1,2\}\}$։ Այդ դեպքում $A \subseteq B$, $B \in C$, և նաև $A \in C$։
\end{itemize}

%======================================================================
\subsection*{Խնդիր 13}
Բերեք բազմությունների օրինակներ, որոնց յուրաքանչյուր տարր նաև այդ բազմության ենթաբազմություն է (տրանզիտիվ բազմություններ)։

$\varnothing$ (դատարկ պայման)։ Նաև $\{\varnothing\}$՝ իր միակ տարրը $\varnothing \subseteq \{\varnothing\}$։
Նաև $\{\varnothing, \{\varnothing\}\}$՝ $\varnothing \subseteq \{\varnothing, \{\varnothing\}\}$ և $\{\varnothing\} \subseteq \{\varnothing, \{\varnothing\}\}$։
Ընդհանուր առմամբ, ցանկացած օրդինալ թիվ տրանզիտիվ բազմություն է։

%======================================================================
\subsection*{Խնդիր 14}
Թվարկեք $2^A$ բազմության տարրերը։

\begin{itemize}
  \item[14.1.] $A = \varnothing$՝ $\boxed{2^A = \{\varnothing\}}$ (1 տարր)։
  \item[14.2.] $A = \{x\}$՝ $\boxed{2^A = \{\varnothing, \{x\}\}}$ (2 տարր)։
  \item[14.3.] $A = \{1,a\}$՝ $\boxed{2^A = \{\varnothing, \{1\}, \{a\}, \{1,a\}\}}$ (4 տարր)։
  \item[14.4.] $A = \{\varnothing, \{1\}, \{2\}\}$՝ $2^A$-ն ունի $2^3 = 8$ տարր.
    \begin{align*}
      2^A = \bigl\{&\varnothing,\;\{\varnothing\},\;\{\{1\}\},\;\{\{2\}\},\\
      &\{\varnothing,\{1\}\},\;\{\varnothing,\{2\}\},\;\{\{1\},\{2\}\},\;\{\varnothing,\{1\},\{2\}\}\bigr\}
    \end{align*}
  \item[14.5.] $A = \{1, \{3\}, \{1,2\}\}$՝ $2^A$-ն ունի $2^3 = 8$ տարր.
    \begin{align*}
      2^A = \bigl\{&\varnothing,\;\{1\},\;\{\{3\}\},\;\{\{1,2\}\},\\
      &\{1,\{3\}\},\;\{1,\{1,2\}\},\;\{\{3\},\{1,2\}\},\;\{1,\{3\},\{1,2\}\}\bigr\}
    \end{align*}
\end{itemize}

%======================================================================
\subsection*{Խնդիր 15}
Ապացուցել $\varnothing \subseteq A \cap B \subseteq A \cup B$։

\textit{Ապացույց.} $\varnothing \subseteq S$ ցանկացած $S$ բազմության համար։ Եթե $x \in A \cap B$, ապա $x \in A$, հետևաբար $x \in A \cup B$։ Ուստի $A \cap B \subseteq A \cup B$։ \qed

%======================================================================
\subsection*{Խնդիր 16}
Նկարագրեք $\overline{A},\ \overline{A\cup B},\ \overline{C},\ A\setminus C,\ C\setminus(A\cup B)$, եթե $U = \mathbb{Z}$, $A = \{2k \mid k \in \mathbb{Z}^+\}$ (դրական զույգ թվեր), $B = \{2k-1 \mid k \in \mathbb{Z}^+\}$ (դրական կենտ թվեր), $C = \{x \in \mathbb{Z} \mid x < 10\}$։

Այստեղ $A \cap B = \varnothing$ և $A \cup B = \mathbb{Z}^+$ (յուրաքանչյուր դրական ամբողջ թիվ զույգ է կամ կենտ)։
\begin{itemize}
  \item $\boxed{\overline{A} = \{x \in \mathbb{Z} \mid x \le 0 \text{ կամ } x \text{-ը կենտ է}\}}$
  \item $\boxed{\overline{A \cup B} = \{x \in \mathbb{Z} \mid x \le 0\} = \{0,-1,-2,\ldots\}}$
  \item $\boxed{\overline{C} = \{x \in \mathbb{Z} \mid x \ge 10\}}$
  \item $\boxed{A \setminus C = \{10, 12, 14, \ldots\}}$ \quad (զույգ թվեր $\ge 10$)
  \item $\boxed{C \setminus (A \cup B) = \{x \in \mathbb{Z} \mid x \le 0\} = \{\ldots, -2, -1, 0\}}$
\end{itemize}

%======================================================================
\subsection*{Խնդիր 17}
$A$ = դրական զույգ թվեր, $B$ = դրական կենտ թվեր, $C$ = 3-ի պատիկ դրական թվեր։

\textbf{(a)}
\begin{itemize}
  \item $A \cap C = \{6, 12, 18, \ldots\}$ = 6-ի պատիկ դրական թվեր։
  \item $B \cup C = \{n \in \mathbb{Z}^+ \mid n \text{-ը կենտ է կամ } 3 \mid n\} = \{1,3,5,6,7,9,11,12,13,\ldots\}$։
  \item $B \setminus C = \{1, 5, 7, 11, 13, 17, \ldots\}$ = 3-ի չբաժանվող դրական կենտ թվեր։
\end{itemize}

\textbf{(b)} Այո, բաշխականության օրենքը՝ $A \cap (B \cup C) = (A \cap B) \cup (A \cap C)$, տեղի ունի բոլոր բազմությունների համար։ Քանի որ $A \cap B = \varnothing$ (չկա դրական թիվ, որը միաժամանակ զույգ է և կենտ), ստանում ենք $A \cap (B \cup C) = \varnothing \cup (A \cap C) = A \cap C$ = 6-ի պատիկ դրական թվեր։ Մյուս կողմից, $A \cap (B \cup C)$-ն բաղկացած է դրական զույգ թվերից, որոնք կենտ են կամ 3-ի պատիկ; «կենտ» մասը անհնար է, ուստի սա 3-ի պատիկ դրական զույգ թվերն են (6-ի պատիկներ)։ Երկու կողմերը համընկնում են։ $\checkmark$

%======================================================================
\subsection*{Խնդիր 18}
Կամայական $A$ բազմության համար՝

$\boxed{A \cap \varnothing = \varnothing}$, \quad $\boxed{A \cup \varnothing = A}$, \quad $\boxed{A \setminus \varnothing = A}$, \quad $\boxed{A \setminus A = \varnothing}$, \quad $\boxed{\varnothing \setminus A = \varnothing}$։

%======================================================================
\subsection*{Խնդիր 19}
Հաշվել՝

\begin{itemize}
  \item $\varnothing \setminus \{\varnothing\} = \boxed{\varnothing}$ \quad (դատարկ բազմությունից հեռացնելու բան չկա)։
  \item $\{\varnothing\} \setminus \{\varnothing\} = \boxed{\varnothing}$ \quad (միակ տարրը՝ $\varnothing$-ն, հեռացվում է)։
  \item $\{\varnothing, \{\varnothing\}\} \setminus \varnothing = \boxed{\{\varnothing, \{\varnothing\}\}}$ \quad (ոչինչ չի հեռացվում)։
  \item $\{\varnothing, \{\varnothing\}\} \setminus \{\varnothing\} = \boxed{\{\{\varnothing\}\}}$ \quad ($\varnothing$ տարրը հեռացվում է; մնում է $\{\varnothing\}$-ն)։
\end{itemize}

%======================================================================
\subsection*{Խնդիր 20}
Ապացուցել համարժեք նկարագրությունները։

\textbf{Մաս 1:} $A \subseteq B \iff \overline{B} \subseteq \overline{A} \iff A \cup B = B \iff A \cap B = A$։

\textit{Ապացույցի ուրվագիծ.} Փակում ենք ցիկլը՝ $A \subseteq B \Rightarrow A \cup B = B \Rightarrow A \cap B = A \Rightarrow \overline{B} \subseteq \overline{A} \Rightarrow A \subseteq B$։
$(A \subseteq B) \Rightarrow (A \cup B = B)$: $A$-ի յուրաքանչյուր տարր արդեն $B$-ում է, ուստի $A \cup B = B$։
$(A \cup B = B) \Rightarrow (A \cap B = A)$: եթե $x \in A$, ապա $x \in A \cup B = B$, ուստի $x \in A \cap B$։
$(A \cap B = A) \Rightarrow (\overline{B} \subseteq \overline{A})$: $A \cap B = A$-ն տալիս է $A \subseteq B$; եթե $x \notin B$, ապա $x \notin A$, այսինքն $\overline{B} \subseteq \overline{A}$։
$(\overline{B} \subseteq \overline{A}) \Rightarrow (A \subseteq B)$: կոնտրապոզիցիա --- եթե $x \in A$, ապա $x \notin \overline{A}$, ուստի $x \notin \overline{B}$, այսինքն $x \in B$։ \qed

\textbf{Մաս 2:} $A \cap B = \varnothing \iff A \subseteq \overline{B} \iff B \subseteq \overline{A}$։

\textit{Ապացույցի ուրվագիծ.} Չհատվող լինելը նշանակում է, որ $A$-ի բոլոր տարրերը $B$-ից դուրս են (այսինքն՝ $\overline{B}$-ում), և սիմետրիկորեն՝ $B$-ի բոլոր տարրերը $A$-ից դուրս են։ \qed

\textbf{Մաս 3:} $A \cup B = U \iff \overline{A} \subseteq B \iff \overline{B} \subseteq A$։

\textit{Ապացույցի ուրվագիծ.} Եթե $A \cup B = U$ և $x \notin A$, ապա $x \in B$; ուստի $\overline{A} \subseteq B$։ Եթե $\overline{A} \subseteq B$ և $x \notin B$, ապա $x \notin \overline{A}$, այսինքն՝ $x \in A$; ուստի $\overline{B} \subseteq A$։ Եթե $\overline{B} \subseteq A$, ապա ցանկացած $x \in U$-ի համար կա՛մ $x \in B$, կա՛մ $x \in \overline{B} \subseteq A$, ուստի $A \cup B = U$։ \qed

%======================================================================
\subsection*{Խնդիր 21}
Ապացուցել $(A \setminus B) \setminus C = A \setminus (B \cup C)$։ Արդյո՞ք $C \subseteq A$ պայմանը համարժեք է։

\textit{Ապացույց.} $x \in (A \setminus B) \setminus C \iff x \in A \land x \notin B \land x \notin C \iff x \in A \land x \notin (B \cup C) \iff x \in A \setminus (B \cup C)$։ \qed

$C \subseteq A$ պայմանը համարժեք \textbf{չէ} այս նույնությանը։ Նույնությունը՝ $(A \setminus B) \setminus C = A \setminus (B \cup C)$, տեղի ունի \emph{բոլոր} $A, B, C$ բազմությունների համար, անկախ նրանից՝ $C \subseteq A$ է թե ոչ։

%======================================================================
\subsection*{Խնդիր 22}
Ապացուցել $(A \setminus B) \setminus C = (A \setminus C) \setminus (B \setminus C)$։

\textit{Ապացույց.}
$x \in (A \setminus C) \setminus (B \setminus C) \iff x \in A \land x \notin C \land \lnot(x \in B \land x \notin C)$։
Քանի որ $x \notin C$, $\lnot(x \in B \land x \notin C)$ պայմանը հանգում է $x \notin B$-ի։
Այսպիսով սա հավասար է $x \in A \land x \notin B \land x \notin C = x \in (A \setminus B) \setminus C$։ \qed

%======================================================================
\subsection*{Խնդիր 23}
Ապացուցել $1 + 2 + \cdots + n = \frac{n(n+1)}{2}$։

\textit{Ապացույց.} Հենք՝ $n=1$: $1 = \frac{1 \cdot 2}{2}$։ Ինդուկցիոն քայլ՝ ենթադրենք ճիշտ է $n=k$-ի համար։ Ապա
$\sum_{i=1}^{k+1} i = \frac{k(k+1)}{2} + (k+1) = \frac{k(k+1)+2(k+1)}{2} = \frac{(k+1)(k+2)}{2}$։ \qed

%======================================================================
\subsection*{Խնդիր 24}
Ապացուցել $1^2 + 2^2 + \cdots + n^2 = \frac{n(n+1)(2n+1)}{6}$։

\textit{Ապացույց.} Հենք՝ $n=1$: $1 = \frac{1 \cdot 2 \cdot 3}{6}$։ Ինդուկցիոն քայլ՝ ենթադրենք $k$-ի համար։ Ապա
$\frac{k(k+1)(2k+1)}{6} + (k+1)^2 = \frac{(k+1)[k(2k+1)+6(k+1)]}{6} = \frac{(k+1)(2k^2+7k+6)}{6} = \frac{(k+1)(k+2)(2k+3)}{6}$։ \qed

%======================================================================
\subsection*{Խնդիր 25}
Ապացուցել $1 + 4 + 7 + \cdots + (3n-2) = \frac{n(3n-1)}{2}$։

\textit{Ապացույց.} Հենք՝ $n=1$: $1 = \frac{1 \cdot 2}{2}$։ Ինդուկցիոն քայլ՝
$\frac{k(3k-1)}{2} + (3(k+1)-2) = \frac{k(3k-1) + 2(3k+1)}{2} = \frac{3k^2+5k+2}{2} = \frac{(k+1)(3k+2)}{2} = \frac{(k+1)(3(k+1)-1)}{2}$։ \qed

%======================================================================
\subsection*{Խնդիր 26}
Ապացուցել $\sum_{i=1}^{n} \frac{1}{(2i-1)(2i+1)} = \frac{n}{2n+1}$։

\textit{Ապացույց.} Հենք՝ $n=1$: $\frac{1}{1 \cdot 3} = \frac{1}{3}$։ Ինդուկցիոն քայլ՝
$\frac{k}{2k+1} + \frac{1}{(2k+1)(2k+3)} = \frac{k(2k+3)+1}{(2k+1)(2k+3)} = \frac{2k^2+3k+1}{(2k+1)(2k+3)} = \frac{(k+1)(2k+1)}{(2k+1)(2k+3)} = \frac{k+1}{2k+3}$։ \qed

%======================================================================
\subsection*{Խնդիր 27}
Ապացուցել $1 + 2 + 2^2 + \cdots + 2^{n-1} = 2^n - 1$։

\textit{Ապացույց.} Հենք՝ $n=1$: $1 = 2^1 - 1$։ Ինդուկցիոն քայլ՝
$(2^k - 1) + 2^k = 2 \cdot 2^k - 1 = 2^{k+1} - 1$։ \qed

%======================================================================
\subsection*{Խնդիր 28}
Ապացուցել $1 + r + r^2 + \cdots + r^{n-1} = \frac{1-r^n}{1-r}$, երբ $r \neq 1$։

\textit{Ապացույց.} Հենք՝ $n=1$: $1 = \frac{1-r}{1-r}$։ Ինդուկցիոն քայլ՝
$\frac{1-r^k}{1-r} + r^k = \frac{1-r^k + r^k(1-r)}{1-r} = \frac{1 - r^{k+1}}{1-r}$։ \qed

%======================================================================
\subsection*{Խնդիր 29}
Ապացուցել $1 + 3 + 5 + \cdots + (2n-1) = n^2$։

\textit{Ապացույց.} Հենք՝ $n=1$: $1 = 1^2$։ Ինդուկցիոն քայլ՝
$k^2 + (2(k+1)-1) = k^2 + 2k + 1 = (k+1)^2$։ \qed

%======================================================================
\subsection*{Խնդիր 30}
Ապացուցել $2 + 4 + 6 + \cdots + 2n = n^2 + n$։

\textit{Ապացույց.} Հենք՝ $n=1$: $2 = 1 + 1$։ Ինդուկցիոն քայլ՝
$(k^2 + k) + 2(k+1) = k^2 + 3k + 2 = (k+1)^2 + (k+1)$։ \qed

%======================================================================
\subsection*{Խնդիր 31}
Ապացուցել ինդուկցիայով՝ $a^{m+n} = a^m \cdot a^n$։

\textit{Ապացույց.} Ֆիքսենք $m$-ը; ինդուկցիա ըստ $n$-ի։ Հենք՝ $n=0$: $a^{m+0} = a^m = a^m \cdot 1 = a^m \cdot a^0$։
Ինդուկցիոն քայլ՝ $a^{m+(k+1)} = a^{(m+k)+1} = a^{m+k} \cdot a = (a^m \cdot a^k) \cdot a = a^m \cdot a^{k+1}$։ \qed

%======================================================================
\subsection*{Խնդիր 32}
Ապացուցել ինդուկցիայով՝ $(ab)^k = a^k b^k$։

\textit{Ապացույց.} Հենք՝ $k=0$: $(ab)^0 = 1 = a^0 b^0$։ (Կամ $k=1$: $(ab)^1 = ab = a^1 b^1$։)
Ինդուկցիոն քայլ՝ $(ab)^{k+1} = (ab)^k \cdot ab = a^k b^k \cdot ab = a^{k+1} b^{k+1}$ (օգտվելով բազմապատկման տեղափոխականությունից)։ \qed

%======================================================================
\subsection*{Խնդիր 33}
Ապացուցել $n^2 > 2n + 1$, երբ $n \ge 3$։

\textit{Ապացույց.} Հենք՝ $n=3$: $9 > 7$։ Ինդուկցիոն քայլ՝ ենթադրենք $k^2 > 2k+1$ ինչ-որ $k \ge 3$-ի համար։
$(k+1)^2 = k^2 + 2k + 1 > (2k+1) + 2k + 1 = 4k + 2 \ge 2(k+1) + 1 + (2k-1) > 2(k+1)+1$, 
քանի որ $2k - 1 \ge 5 > 0$։ \qed

%======================================================================
\subsection*{Խնդիր 34}
Ապացուցել $2^n > n^2$, երբ $n \ge 5$։

\textit{Ապացույց.} Հենք՝ $n=5$: $32 > 25$։ Ինդուկցիոն քայլ՝ ենթադրենք $2^k > k^2$, երբ $k \ge 5$։
$2^{k+1} = 2 \cdot 2^k > 2k^2$։ Մեզ պետք է $2k^2 \ge (k+1)^2 = k^2 + 2k + 1$, այսինքն՝ $k^2 - 2k - 1 \ge 0$, որը ճիշտ է $k \ge 3$-ի համար։ \qed

%======================================================================
\subsection*{Խնդիր 35}
Ապացուցել $a^3 > 3a^2 + 3a + 1$, երբ $a \ge 4$։

\textit{Ապացույց.} Հենք՝ $a=4$: $64 > 48 + 12 + 1 = 61$։ $\checkmark$

Ինդուկցիոն քայլ՝ ենթադրենք $k^3 > 3k^2 + 3k + 1$, երբ $k \ge 4$։ Ցույց տանք, որ $(k+1)^3 > 3(k+1)^2 + 3(k+1) + 1$։
$(k+1)^3 = k^3 + 3k^2 + 3k + 1 > (3k^2 + 3k + 1) + 3k^2 + 3k + 1 = 6k^2 + 6k + 2$։
Մեզ պետք է $6k^2 + 6k + 2 \ge 3k^2 + 9k + 7$, այսինքն՝ $3k^2 - 3k - 5 \ge 0$, որը ճիշտ է $k \ge 2$-ի համար։ \qed

%======================================================================
\subsection*{Խնդիր 36}
Ապացուցել, որ $5^{2n} + 3 \cdot 2^{5n-2}$-ը բաժանվում է 7-ի ցանկացած $n$ դրական ամբողջ թվի համար։

\textit{Ապացույց.} Հենք՝ $n=1$: $25 + 3 \cdot 8 = 25 + 24 = 49 = 7 \cdot 7$։ $\checkmark$

Ինդուկցիոն քայլ՝ ենթադրենք $7 \mid (5^{2k} + 3 \cdot 2^{5k-2})$։ Ապա
$5^{2(k+1)} + 3 \cdot 2^{5(k+1)-2} = 25 \cdot 5^{2k} + 32 \cdot 3 \cdot 2^{5k-2}$
$= 25 \cdot 5^{2k} + 25 \cdot 3 \cdot 2^{5k-2} + 7 \cdot 3 \cdot 2^{5k-2}$
$= 25(5^{2k} + 3 \cdot 2^{5k-2}) + 7 \cdot 3 \cdot 2^{5k-2}$։
Երկու գումարելիներն էլ բաժանվում են 7-ի։ \qed

%======================================================================
\subsection*{Խնդիր 37}
Ապացուցել, որ $3^{n+2} + 4^{2n+1}$-ը բաժանվում է 13-ի ցանկացած $n \ge 0$-ի համար։

\textit{Ապացույց.} Հենք՝ $n=0$: $9 + 4 = 13$։ $\checkmark$ (Կամ $n=1$: $27 + 64 = 91 = 7 \cdot 13$։)

Ինդուկցիոն քայլ՝ ենթադրենք $13 \mid (3^{k+2} + 4^{2k+1})$։ Ապա
$3^{k+3} + 4^{2k+3} = 3 \cdot 3^{k+2} + 16 \cdot 4^{2k+1}$
$= 3(3^{k+2} + 4^{2k+1}) + 13 \cdot 4^{2k+1}$։
Երկու գումարելիներն էլ բաժանվում են 13-ի։ \qed

%======================================================================
\subsection*{Խնդիր 38}
Ապացուցել, որ $5^{2n} + 7$-ը բաժանվում է 8-ի ցանկացած $n$ դրական ամբողջ թվի համար։

\textit{Ապացույց.} $5^{2n} = 25^n \equiv 1^n = 1 \pmod{8}$, ուստի $5^{2n} + 7 \equiv 1 + 7 = 8 \equiv 0 \pmod{8}$։

Այլ տարբերակ ինդուկցիայով՝ Հենք՝ $n=1$: $25 + 7 = 32 = 8 \cdot 4$։
Ինդուկցիոն քայլ՝ $5^{2(k+1)} + 7 = 25 \cdot 5^{2k} + 7 = 25(5^{2k}+7) - 25 \cdot 7 + 7 = 25(5^{2k}+7) - 168$։
Ե՛վ $25(5^{2k}+7)$-ը, և՛ $168 = 8 \cdot 21$-ը բաժանվում են 8-ի։ \qed

%======================================================================
\subsection*{Խնդիր 39}
Ապացուցել, որ $3^{3n+1} + 2^{n+1}$-ը բաժանվում է 5-ի ($n \in \mathbb{N}$)։

\textit{Ապացույց.} Հենք՝ $n=0$: $3 + 2 = 5$։ $\checkmark$

Ինդուկցիոն քայլ՝ $3^{3(k+1)+1} + 2^{k+2} = 27 \cdot 3^{3k+1} + 2 \cdot 2^{k+1}$
$= 27 \cdot 3^{3k+1} + 27 \cdot 2^{k+1} - 25 \cdot 2^{k+1}$
$= 27(3^{3k+1} + 2^{k+1}) - 25 \cdot 2^{k+1}$։
Ըստ ենթադրության $5 \mid 27(3^{3k+1} + 2^{k+1})$ (քանի որ $5$-ը բաժանում է փակագծի արտահայտությունը) և $5 \mid 25 \cdot 2^{k+1}$։ \qed

%======================================================================
\subsection*{Խնդիր 40}
Ապացուցել, որ $9^n - 1$-ը բաժանվում է 8-ի ($n \in \mathbb{N}$)։

\textit{Ապացույց.} $9^n - 1 = (9-1)(9^{n-1} + 9^{n-2} + \cdots + 1) = 8 \cdot \sum_{i=0}^{n-1} 9^i$։

Այլ տարբերակ՝ $9 \equiv 1 \pmod{8}$, ուստի $9^n \equiv 1 \pmod{8}$։ \qed

%======================================================================
\subsection*{Խնդիր 41}
Ապացուցել, որ $10^n - 1$-ը բաժանվում է 9-ի ($n \in \mathbb{N}$)։

\textit{Ապացույց.} $10 \equiv 1 \pmod{9}$, ուստի $10^n \equiv 1 \pmod{9}$, հետևաբար $9 \mid (10^n - 1)$։

Այլ տարբերակ՝ $10^n - 1 = \underbrace{99\ldots9}_{n} = 9 \cdot \underbrace{11\ldots1}_{n}$։ \qed

%======================================================================
\subsection*{Խնդիր 42}
Ապացուցել, որ $n^2 - 1$-ը բաժանվում է 4-ի բոլոր կենտ բնական թվերի համար։

\textit{Ապացույց.} Դիցուք $n = 2m+1$, որտեղ $m \ge 0$։ Ապա $n^2 - 1 = (2m+1)^2 - 1 = 4m^2 + 4m = 4m(m+1)$, որը բաժանվում է 4-ի։ \qed

(Ինդուկցիայի կարիք չկա; ուղղակի ապացույցը բավարար է։)

%======================================================================
\subsection*{Խնդիր 43}
Գտնել և ապացուցել՝ $3 + 7 + 11 + \cdots + (4n-1)$։

Գումարն է՝ $\boxed{n(2n+1)}$։

\textit{Ապացույց.} Հենք՝ $n=1$: $3 = 1 \cdot 3$։ $\checkmark$

Ինդուկցիոն քայլ՝ $k(2k+1) + (4(k+1)-1) = 2k^2 + k + 4k + 3 = 2k^2 + 5k + 3 = (k+1)(2k+3) = (k+1)(2(k+1)+1)$։ \qed

%======================================================================
\subsection*{Խնդիր 44}
Գտնել և ապացուցել՝ $1 + 5 + 9 + \cdots + (4n+1)$։

Վերջին $4n+1$ անդամը $1,5,9,\ldots$ հաջորդականության $(n+1)$-րդ անդամն է, ուստի գումարն ունի $n+1$ անդամ և հավասար է $\boxed{(n+1)(2n+1)}$-ի։

\textit{Ապացույց.} $\displaystyle\sum_{k=1}^{n+1}(4k-3) = 4\cdot\frac{(n+1)(n+2)}{2} - 3(n+1) = (n+1)\bigl[2(n+2)-3\bigr] = (n+1)(2n+1)$։
Ստուգում՝ $n=1$: $1+5 = 6 = 2\cdot 3$;\; $n=2$: $1+5+9 = 15 = 3\cdot 5$։ $\checkmark$ \qed

%======================================================================
\subsection*{Խնդիր 45}
Գտնել և ապացուցել՝ $-1 + 2 - 3 + 4 - \cdots - (2n-1) + 2n$։

Գումարն է՝ $\boxed{n}$։

\textit{Ապացույց.} Խմբավորենք $n$ զույգով՝ $\bigl(-(2k-1) + 2k\bigr) = 1$ $k = 1, \ldots, n$-ի համար։ Ուստի ընդհանուրը $n \cdot 1 = n$ է։
Ստուգում՝ $-1+2 = 1$;\; $-1+2-3+4 = 2$;\; $-1+2-3+4-5+6 = 3$։ $\checkmark$ \qed

%======================================================================
\subsection*{Խնդիր 46}
Գտնել և ապացուցել՝ $1 - 3 + 5 - 7 + \cdots + (4n-3) - (4n-1)$։

Գումարն է՝ $\boxed{-2n}$։

\textit{Ապացույց.} Խմբավորենք $n$ զույգով՝ $\bigl((4k-3) - (4k-1)\bigr) = -2$ $k = 1, \ldots, n$-ի համար։ Ուստի ընդհանուրը $-2n$ է։
Ստուգում՝ $1-3 = -2$;\; $1-3+5-7 = -4$;\; $1-3+5-7+9-11 = -6$։ $\checkmark$ \qed

%======================================================================
\subsection*{Խնդիր 47}
Ապացուցել $2^n > n^3$ և $n! > 100^n$ մեծ $n$-երի համար։

\textbf{(a)} $2^n > n^3$, երբ $n \ge 10$։

\textit{Ապացույց.} Հենք՝ $n=10$: $1024 > 1000$։ Ինդուկցիոն քայլ՝ ենթադրենք $2^k > k^3$, երբ $k \ge 10$։ Ապա $2^{k+1} = 2 \cdot 2^k > 2k^3$։ Մեզ պետք է $2k^3 \ge (k+1)^3$, այսինքն՝ $(1+1/k)^3 \le 2$։ $k \ge 10$ դեպքում՝ $(1+1/k)^3 \le (1.1)^3 = 1.331 < 2$։ $\checkmark$ \qed

\textbf{(b)} $n! > 100^n$, երբ $n \ge 200$ (օրինակ)։

\textit{Ապացույց.} Ինչ-որ պահից սկսած, երբ $k \ge 200$, ունենք $k + 1 > 100$, ուստի $(k+1)! = (k+1) \cdot k! > 100 \cdot 100^k = 100^{k+1}$. Հենքը՝ $200! > 100^{200}$ կարելի է ստուգել Ստիռլինգի բանաձևով։ \qed

%======================================================================
\subsection*{Խնդիր 48}
Հանոյի աշտարակ. ապացուցել, որ $n$ սկավառակը կարելի է տեղափոխել $2^n - 1$ քայլով։

\textit{Ապացույց.} Հենք՝ $n=1$: տեղափոխել միակ սկավառակը ուղղակիորեն (1 քայլ $= 2^1 - 1$)։
Ինդուկցիոն քայլ՝ $k+1$ սկավառակ $A$-ից $C$ տեղափոխելու համար նախ տեղափոխում ենք վերևի $k$ սկավառակները $A$-ից $B$ (ըստ վարկածի՝ $2^k - 1$ քայլ), հետո ամենամեծ սկավառակը $A$-ից $C$ (1 քայլ), այնուհետև $k$ սկավառակները $B$-ից $C$ ($2^k - 1$ քայլ)։ Ընդհանուր՝ $2(2^k-1)+1 = 2^{k+1}-1$։ \qed

%======================================================================
\subsection*{Խնդիր 49}
Բեռնուլիի անհավասարություն. $(1+h)^n \ge 1 + nh$, երբ $h > -1$, $n \in \mathbb{N}$։

\textit{Ապացույց.} Հենք՝ $n=1$: $1+h \ge 1+h$։ $\checkmark$

Ինդուկցիոն քայլ՝ $(1+h)^{k+1} = (1+h)^k(1+h) \ge (1+kh)(1+h) = 1 + (k+1)h + kh^2 \ge 1 + (k+1)h$,
քանի որ $kh^2 \ge 0$։ (Օգտվեցինք նրանից, որ $1+h > 0$, անհավասարության նշանը պահպանելու համար)։ \qed

%======================================================================
\subsection*{Խնդիր 50}
Ապացուցել Թեորեմ 1-ը (բազմությունների հանրահաշվի հիմնական օրենքները)՝ օգտվելով պատկանելության հարաբերությունից։

«Թեորեմ 1»-ը սովորաբար վերաբերում է հիմնական նույնություններին՝
\begin{enumerate}
  \item Տեղափոխականություն՝ $A \cup B = B \cup A$, $A \cap B = B \cap A$։
  \item Զուգորդականություն՝ $(A \cup B) \cup C = A \cup (B \cup C)$, նմանապես $\cap$-ի համար։
  \item Բաշխականություն՝ $A \cup (B \cap C) = (A \cup B) \cap (A \cup C)$, $A \cap (B \cup C) = (A \cap B) \cup (A \cap C)$։
  \item Նույնություն՝ $A \cup \varnothing = A$, $A \cap U = A$։
  \item Լրացում՝ $A \cup \overline{A} = U$, $A \cap \overline{A} = \varnothing$։
  \item Իդեմպոտենտություն՝ $A \cup A = A$, $A \cap A = A$։
  \item Դե Մորգան՝ $\overline{A \cup B} = \overline{A} \cap \overline{B}$, $\overline{A \cap B} = \overline{A} \cup \overline{B}$։
\end{enumerate}

\textit{Ապացույցի ուրվագիծ.} Յուրաքանչյուր նույնություն ապացուցվում է ցույց տալով, որ $x \in \text{Ձախ մաս} \iff x \in \text{Աջ մաս}$ կամայական $x$-ի համար, թարգմանելով բազմային գործողությունները տրամաբանական գործողությունների ($\cup \to \lor$, $\cap \to \land$, $\overline{\cdot} \to \lnot$) և կիրառելով ստանդարտ տրամաբանական օրենքները։ \qed

%======================================================================
\subsection*{Խնդիր 51}
Ապացուցել Թեորեմ 2-ը՝ օգտվելով պատկանելության հարաբերությունից։

«Թեորեմ 2»-ը սովորաբար ներառում է լրացուցիչ նույնություններ՝ կլանում ($A \cup (A \cap B) = A$, $A \cap (A \cup B) = A$), ինվոլյուցիա ($\overline{\overline{A}} = A$), գերակայություն ($A \cup U = U$, $A \cap \varnothing = \varnothing$), և $\overline{U} = \varnothing$, $\overline{\varnothing} = U$։

\textit{Ապացույցի ուրվագիծ.} Ինչպես Խնդիր 50-ում, յուրաքանչյուրն ապացուցվում է տարր առ տարր՝ վերածվելով ասույթային տրամաբանության նույնաբանությունների։ Օրինակ՝ կլանում՝ $x \in A \cup (A \cap B) \iff x \in A \lor (x \in A \land x \in B) \iff x \in A$։ \qed

%======================================================================
\subsection*{Խնդիր 52}
Ապացուցել Թեորեմ 2-ը՝ օգտվելով Թեորեմ 1-ի նույնություններից որպես աքսիոմներ (հանրահաշվական ապացույց, առանց տարր առ տարր)։

\textit{Ապացույցի ուրվագիծ.} Օրինակ՝ կլանում՝ $A \cup (A \cap B) = (A \cap U) \cup (A \cap B) = A \cap (U \cup B) = A \cap U = A$, օգտվելով Թեորեմ 1-ի նույնության, բաշխականության և գերակայության օրենքներից։ Թեորեմ 2-ի յուրաքանչյուր նույնություն կարելի է ստանալ զուտ հանրահաշվորեն Թեորեմ 1-ից։ \qed

%======================================================================
\subsection*{Խնդիր 53}
Թեորեմ 1-ի և 2-ի օգնությամբ ապացուցել նույնությունները։

\textbf{(1)} $(A \cap B \cap X) \cup (A \cap B \cap C \cap X \cap Y) \cup (A \cap X \cap \overline{A}) = A \cap B \cap X$.

\textit{Ապացույցի ուրվագիծ.} Երրորդ անդամը պարունակում է $A \cap \overline{A} = \varnothing$, ուստի այն անհետանում է։ Երկրորդ անդամն առաջինի ենթաբազմություն է (ունի լրացուցիչ $C$ և $Y$ պայմաններ), ուստի առաջին $\cup$ երկրորդ $=$ առաջին՝ կլանմամբ։ Արդյունք՝ $A \cap B \cap X$։ \qed

\textbf{(2)} $(A \cap B \cap C) \cup (\overline{A} \cap B \cap C) \cup \overline{B} \cup \overline{C} = U$.

\textit{Ապացույցի ուրվագիծ.} $(A \cap B \cap C) \cup (\overline{A} \cap B \cap C) = (A \cup \overline{A}) \cap B \cap C = B \cap C$, ուստի ձախ կողմը $(B \cap C) \cup \overline{B} \cup \overline{C} = (B \cap C) \cup \overline{B \cap C} = U$ է (Դե Մորգան)։ \qed

\textbf{(3)} $(A \cap B \cap C \cap \overline{X}) \cup (\overline{A} \cap C) \cup (\overline{B} \cap C) \cup (C \cap X) = C$.

\textit{Ապացույցի ուրվագիծ.} Դուրս բերենք $C$-ն՝ ձախ կողմը $C \cap \bigl[(A \cap B \cap \overline{X}) \cup \overline{A} \cup \overline{B} \cup X\bigr]$ է։ $S = A \cap B \cap \overline{X}$-ով՝ $\overline{A} \cup \overline{B} \cup X = \overline{S}$, ուստի փակագիծը $S \cup \overline{S} = U$ է։ Հետևաբար արդյունքը $C \cap U = C$ է։ \qed

%======================================================================
\subsection*{Խնդիր 54}
Ապացուցել բազմությունների նույնությունները (բոլորը ճիշտ; յուրաքանչյուրը՝ տարր առ տարր, $\setminus$-ը գրելով որպես $\cap\,\overline{(\cdot)}$, իսկ $A \triangle B$-ն՝ «$x$-ը $A, B$-ից ճիշտ մեկում»)։

\textbf{(1)} $A \cap (B \setminus C) = (A \cap B) \setminus C$. \\
Երկու կողմում էլ՝ $x \in A \land x \in B \land x \notin C$։ \checkmark

\textbf{(2)} $A \setminus (B \cup C) = (A \setminus B) \setminus C$. \\
Երկու կողմում էլ՝ $x \in A \land x \notin B \land x \notin C$։ \checkmark

\textbf{(3)} $(A \cup B) \setminus C = (A \setminus C) \cup (B \setminus C)$. \\
$(x \in A \lor x \in B) \land x \notin C \iff (x \in A \land x \notin C) \lor (x \in B \land x \notin C)$։ \checkmark

\textbf{(4)} $A \cap (B \setminus C) = (A \cap B) \setminus (A \cap C)$. \\
$x \in A$-ի դեպքում $\lnot(x \in A \land x \in C)$-ն վերածվում է $x \notin C$-ի; երկու կողմն էլ՝ $x \in A \land x \in B \land x \notin C$։ \checkmark

\textbf{(5)} $A \setminus (B \setminus C) = (A \setminus B) \cup (A \cap C)$. \\
$x \in A \land (x \notin B \lor x \in C) \iff (x \in A \land x \notin B) \lor (x \in A \land x \in C)$։ \checkmark

\textbf{(6)} $(A \setminus B) \setminus C = (A \setminus C) \setminus (B \setminus C)$. \\
$x \notin C$-ի դեպքում $B \setminus C$-ն վերածվում է $B$-ի; երկու կողմն էլ՝ $x \in A \land x \notin B \land x \notin C$ (հմմտ.\ Խնդիր 22)։ \checkmark

\textbf{(7)} $A \triangle \varnothing = A$. \quad $(A \setminus \varnothing) \cup (\varnothing \setminus A) = A \cup \varnothing = A$։ \checkmark

\textbf{(8)} $A \triangle B = B \triangle A$. \quad «$A, B$-ից ճիշտ մեկում»-ը սիմետրիկ է $A, B$-ի նկատմամբ։ \checkmark

\textbf{(9)} $A \triangle (B \triangle C) = (A \triangle B) \triangle C$. \quad Երկու կողմն էլ ճիշտ են, երբ $x$-ը պատկանում է $A, B, C$-ից \emph{կենտ} թվով բազմության (զույգությունն ասոցիատիվ է)։ \checkmark

\textbf{(10)} $A \triangle A = \varnothing$. \quad $(A \setminus A) \cup (A \setminus A) = \varnothing$։ \checkmark

\textbf{(11)} $A \cap (B \triangle C) = (A \cap B) \triangle (A \cap C)$. \\
Երկու կողմն էլ վերածվում են $(x \in A \land x \in B \land x \notin C) \lor (x \in A \land x \notin B \land x \in C)$-ի։ \checkmark

%======================================================================
\subsection*{Խնդիր 55}
Ցույց տալ, որ $A \cup (B \setminus C) \neq (A \cup B) \setminus (A \cup C)$ ընդհանուր դեպքում։

Հակաօրինակ՝ $A = \{1\}$, $B = \{1,2\}$, $C = \{1\}$։
Ձախ՝ $B \setminus C = \{2\}$, ուստի $A \cup (B \setminus C) = \{1,2\}$։
Աջ՝ $A \cup B = \{1,2\}$, $A \cup C = \{1\}$, ուստի $(A \cup B) \setminus (A \cup C) = \{2\}$։
$\{1,2\} \neq \{2\}$։ \qed

%======================================================================
\subsection*{Խնդիր 56}
Ապացուցել, որ $A \triangle X = B$-ն ունի միակ լուծում։

Միակ լուծումն է՝ $\boxed{X = A \triangle B}$։

\textit{Ապացույց.}
\textbf{Գոյություն:} Դիցուք $X = A \triangle B$։ Ապա $A \triangle X = A \triangle (A \triangle B) = (A \triangle A) \triangle B = \varnothing \triangle B = B$։ $\checkmark$

\textbf{Միակություն:} Եթե $A \triangle X = B$, երկու կողմին կիրառենք $A \triangle$՝ $A \triangle (A \triangle X) = A \triangle B$, որից $X = A \triangle B$։ (Օգտվեցինք $\triangle$-ի ասոցիատիվությունից և $A \triangle A = \varnothing$, $\varnothing \triangle X = X$-ից։) \qed

%======================================================================
\subsection*{Խնդիր 57}
Գոյություն ունե՞ն $A, B$ բազմություններ, որ $|A \times B|$ հավասար է՝

\textbf{57.1.} 12 տարր։ \textbf{ԱՅՈ։} Օրինակ՝ $A = \{1,2,3\}$, $B = \{a,b,c,d\}$; $|A| \cdot |B| = 3 \cdot 4 = 12$։

\textbf{57.2.} 13 տարր։ \textbf{ԱՅՈ։} Քանի որ $13$-ը պարզ է, $|A| \cdot |B| = 13$ պահանջում է $|A|=1, |B|=13$ կամ հակառակը։ Օրինակ՝ $A=\{1\}$, $B$՝ ցանկացած 13-տարրանոց բազմություն։

$\boxed{\text{57.1: Այո (օր.\ } 3 \times 4 = 12\text{)։ \quad 57.2: Այո (օր.\ } 1 \times 13 = 13\text{)։}}$

%======================================================================
\subsection*{Խնդիր 58}
Ապացուցել, որ $A \times B = B \times A \implies A = B$ (ենթադրելով երկուսն էլ ոչ դատարկ)։

\textit{Ապացույց.} Ենթադրենք $A, B \neq \varnothing$ և $A \times B = B \times A$։
Ցույց տալու համար $A \subseteq B$՝ դիցուք $a \in A$, ընտրենք ցանկացած $b \in B$։ Ապա $(a,b) \in A \times B = B \times A$, ուստի $a \in B$։
Ցույց տալու համար $B \subseteq A$՝ դիցուք $b \in B$, ընտրենք ցանկացած $a \in A$։ Ապա $(b,a) \in B \times A = A \times B$, ուստի $b \in A$։
Հետևաբար $A = B$։ (Եթե որևէ բազմություն դատարկ է, երկու արտադրյալներն էլ $\varnothing$ են՝ անկախ $A = B$-ից։) \qed

%======================================================================
\subsection*{Խնդիր 59}
Թվարկել $A \times B$-ի տարրերը՝ $A = \{a,b\}$, $B = \{1,2,3,4\}$-ի համար։

\[
  \boxed{A \times B = \{(a,1),(a,2),(a,3),(a,4),(b,1),(b,2),(b,3),(b,4)\}}
\]
(8 տարր։)

%======================================================================
\subsection*{Խնդիր 60}
Ապացուցել, որ դեկարտյան արտադրյալը բաշխվում է բազմությունների գործողությունների վրա։

\textbf{(1)} $A \times (B \cup C) = (A \times B) \cup (A \times C)$.

\textit{Ապացույց.} $(x,y) \in A \times (B \cup C) \iff x \in A \land (y \in B \lor y \in C) \iff (x \in A \land y \in B) \lor (x \in A \land y \in C) \iff (x,y) \in (A\times B) \cup (A \times C)$։ \qed

\textbf{(2)} $(B \cup C) \times A = (B \times A) \cup (C \times A)$. Նույն ապացույցը՝ առաջին բաղադրիչի վրա։

\textbf{(3)} $A \times (B \cap C) = (A \times B) \cap (A \times C)$.

\textit{Ապացույց.} $(x,y) \in$ ձախ $\iff x \in A \land y \in B \land y \in C \iff (x \in A \land y \in B) \land (x \in A \land y \in C)$։ \qed

\textbf{(4)} $(B \cap C) \times A = (B \times A) \cap (C \times A)$. (3)-ի հայելին։

\textbf{(5)} $A \times (B \setminus C) = (A \times B) \setminus (A \times C)$.

\textit{Ապացույց.} $(x,y) \in$ ձախ $\iff x \in A \land y \in B \land y \notin C$։ $(x,y) \in$ աջ $\iff (x \in A \land y \in B) \land \lnot(x \in A \land y \in C)$։ $x \in A$-ի դեպքում երկրորդ պայմանը վերածվում է $y \notin C$-ի։ \qed

\textbf{(6)} $(B \setminus C) \times A = (B \times A) \setminus (C \times A)$. (5)-ի հայելին։

%======================================================================
\subsection*{Խնդիր 61}
Ապացուցել $(A \cap B) \times (C \cap D) = (A \times C) \cap (B \times D)$։

\textit{Ապացույց.}
$(x,y) \in (A \cap B) \times (C \cap D)$
$\iff x \in A \land x \in B \land y \in C \land y \in D$
$\iff (x \in A \land y \in C) \land (x \in B \land y \in D)$
$\iff (x,y) \in (A \times C) \cap (B \times D)$։ \qed

%======================================================================
\subsection*{Խնդիր 62}
Ապացուցել $A \subseteq B \land C \subseteq D \implies A \times C \subseteq B \times D$։

\textit{Ապացույց.} Դիցուք $(x,y) \in A \times C$։ Ապա $x \in A \subseteq B$ տալիս է $x \in B$, իսկ $y \in C \subseteq D$ տալիս է $y \in D$։ Ուստի $(x,y) \in B \times D$։ \qed